% Standalone supplementary-methods document for sans2d-mode-analyzer.
% Layout reference: ../../Michels_et_al_2026_vortexNPs_SM.tex.
% Compile twice with pdflatex; no external figures or bibliography are required.
\documentclass[aps,prb,preprint,onecolumn,notitlepage,amsmath,amssymb,
  superscriptaddress,floatfix]{revtex4-2}
\usepackage[T1]{fontenc}
\usepackage[utf8]{inputenc}
\usepackage{bm}
\usepackage{graphicx}
\usepackage{etoolbox}
\usepackage{hyperref}
\hypersetup{colorlinks=true,linkcolor=blue,urlcolor=blue,
  pdftitle={Supplemental Material: Orthogonal mode analysis of two-dimensional small-angle neutron scattering},
  pdfauthor={SANS2D mode analyzer contributors}}

% Retain the preprint typography without its forced break before the abstract.
\makeatletter
\patchcmd{\frontmatter@abstract@produce}{\preprintsty@sw}{\@secondoftwo}
  {}{\ClassError{revtex4-2}{Cannot adjust abstract spacing}{}}
\makeatother

\newcommand{\dd}{\,\mathrm{d}}
\newcommand{\Ann}{\mathcal{A}}
\newcommand{\Modes}{\mathcal{K}}
\newcommand{\trans}{\mathsf{T}}
\newcommand{\tr}{\operatorname{tr}}
\renewcommand{\theequation}{S\arabic{equation}}
\renewcommand{\thesection}{S\arabic{section}}


\begin{document}
\raggedbottom

\title{Supplemental Material:
Orthogonal mode analysis of two-dimensional small-angle neutron scattering}
\author{SANS2D mode analyzer contributors}
\noaffiliation
\date{13 November 2026}

\begin{abstract}
This document describes the analysis implemented in
\texttt{sans2d-mode-analyzer} (version 0.1.0). A real-valued detector pattern
on a circular annulus is represented by orthonormal radial Legendre modes
and real angular Fourier modes. We specify the linear fit, its covariance,
the associated Parseval powers, angular block statistics and optional numerical
Holm diagnostics, and the radial coefficient spectra. We also describe detector selection, optional pixel
integration and parity-based centering, and symmetry diagnostics.
All quantities refer to the chosen annulus and finite fitted expansion. We
finally distinguish numerical verification, data-level validation, and
physical interpretation as separate stages of the analysis.
\end{abstract}

\maketitle

\tableofcontents
\clearpage

\section{Methodology of the two-dimensional SANS mode analysis}
\label{sec:methodology}

\subsection{Scope, detector coordinates, and integration measure}

The input is a reduced, real-valued scattering intensity or differential
cross section, denoted by $I(q,\theta)$. The analysis describes its radial
and angular dependence without prescribing a particle form factor or a
magnetization model. In particular, the modal fit does not estimate vortex
amplitudes, particle radii, or structure-factor parameters. Such physical
interpretations require a separate model comparison.

The detector coordinates are $q_z$ along the chosen reference axis and
$q_y$ transverse to it. For the magnetic-SANS geometry of the accompanying
nanoflower manuscript, the reference axis is the applied-field direction.
The code uses
\begin{equation}
 q=\sqrt{q_z^2+q_y^2},\qquad
 \theta=\operatorname{atan2}(q_y,q_z),\qquad
 (q_z,q_y)=q(\cos\theta,\sin\theta).
 \label{eq:coordinates}
\end{equation}
Angles are understood modulo $2\pi$. The analysis domain and its area are
\begin{equation}
 \Ann=\{(q,\theta):q_{\min}\leq q\leq q_{\max},\ 0\leq\theta<2\pi\},
 \qquad A=\pi\Delta,\qquad
 \Delta=q_{\max}^2-q_{\min}^2>0.
 \label{eq:annulus}
\end{equation}
Here $q_{\min}\geq0$; when the origin is included, its angular coordinate
has no effect on continuous integrals because it has zero area.

\begin{samepage}
For a radial-measure exponent $s>-1$, the inner product is
\begin{equation}
 \langle f,g\rangle_{\Ann,s}
 =\int_{q_{\min}}^{q_{\max}}\int_0^{2\pi}
 f(q,\theta)g(q,\theta)\dd\theta\,q^s\dd q.
 \label{eq:inner-product}
\end{equation}
\end{samepage}
The default $s=1$ is the two-dimensional detector-plane area measure
$q\dd q\dd\theta$. Other exponents intentionally reweight the radial
interval; for example, $s>1$ emphasizes larger $q$. The choice is not a
change of the experimental cross section and should be reported with the
analysis.

\subsection{Orthonormal expansion and mode indices}

The radial coordinate and radial modes are
\begin{align}
 \Delta_s&=q_{\max}^{s+1}-q_{\min}^{s+1},\nonumber\\
 x_s(q)&=\frac{2q^{s+1}-q_{\max}^{s+1}-q_{\min}^{s+1}}{\Delta_s},
 \qquad q^s\dd q=\frac{\Delta_s}{2(s+1)}\dd x_s,
 \label{eq:radial-coordinate}\\
 R_m^{(s)}(q)&=\sqrt{\frac{(2m+1)(s+1)}{\Delta_s}}\,P_m[x_s(q)],
 \qquad m=0,1,\ldots,
 \label{eq:radial-modes}
\end{align}
where $P_m$ is the Legendre polynomial of degree $m$. Since
$\int_{-1}^{1}P_m(x)P_l(x)\dd x=2\delta_{ml}/(2m+1)$,
\begin{equation}
 \int_{q_{\min}}^{q_{\max}}R_m^{(s)}(q)R_l^{(s)}(q)\,q^s\dd q=\delta_{ml}.
 \label{eq:radial-orthogonality}
\end{equation}
At the default $s=1$, these expressions reduce exactly to the coordinate
affine in $q^2$ and the normalization used previously.

To distinguish the angular basis from the block statistic $T_n$ below,
we write the angular functions as $\mathcal{T}_n$:
\begin{equation}
 \mathcal{T}_n(\theta)=
 \begin{cases}
 (2\pi)^{-1/2},&n=0,\\
 \pi^{-1/2}\cos(n\theta),&n>0,\\
 \pi^{-1/2}\sin(|n|\theta),&n<0.
 \end{cases}
 \qquad
 \int_0^{2\pi}\mathcal{T}_n\mathcal{T}_{n'}\dd\theta=\delta_{nn'}.
 \label{eq:angular-modes}
\end{equation}
Thus the signed index $n$ encodes the trigonometric branch; it is not a
negative angular frequency of a complex Fourier expansion. The product modes
and fitted intensity are
\begin{equation}
 \Phi_{nm}^{(s)}(q,\theta)=R_m^{(s)}(q)\mathcal{T}_n(\theta),\qquad
 \widehat I(q,\theta)=\sum_{(n,m)\in\Modes}
 \widehat c_{nm}\Phi_{nm}^{(s)}(q,\theta),
 \label{eq:expansion}
\end{equation}
with $\langle\Phi_{nm}^{(s)},\Phi_{n'm'}^{(s)}\rangle_{\Ann,s}
=\delta_{nn'}\delta_{mm'}$. By default, the standard analyzer uses the full
rectangle $|n|\leq n_{\max}$ and $0\leq m\leq m_{\max}$, containing
$K=(2n_{\max}+1)(m_{\max}+1)$ product modes. With \texttt{even\_only=True},
the fit uses the subset with even $|n|$ at every radial order, so that
$\widehat I(q,\theta+\pi)=\widehat I(q,\theta)$. The lower-level fitting
functions accept any explicitly selected set $\Modes$. Omitted coefficients
are understood as zero in subsequent sums.

For a continuous, unweighted orthogonal projection of a square-integrable
function, the coefficients would be
\begin{equation}
 c_{nm}^{\mathrm{proj}}=\langle I,\Phi_{nm}^{(s)}\rangle_{\Ann,s}.
 \label{eq:projection}
\end{equation}
The implemented fit instead estimates them jointly from the available
detector samples. Masks, nonuniform sampling, experimental uncertainties,
and finite pixels generally destroy orthogonality of the sampled columns.
Consequently, Eq.~\eqref{eq:projection} is not used as a shortcut for fitting
masked or uncertainty-weighted detector data.

\subsection{Sample selection, linear fitting, and covariance}

Let $p=1,\ldots,N$ enumerate the admitted detector samples and $j$ enumerate
the selected pairs $(n,m)$. In the default point-evaluation approximation,
$X_{pj}=\Phi_j(q_p,\theta_p)$ and $y_p=I_p$. Selection requires finite
coordinates and intensity, and a sample center in the chosen annulus.
If uncertainties $\sigma_p$ or sample weights $w_p$ are supplied, their
selected values must additionally be finite and strictly positive.
Finite zero and negative background-corrected intensities remain in the
fit. Logarithmic color scales affect display only.

The objective is
\begin{equation}
 Q(\bm c)=(\bm y-X\bm c)^{\trans}W(\bm y-X\bm c),\qquad
 W_{pp}=\frac{w_p}{\sigma_p^2}.
 \label{eq:wls}
\end{equation}
Absent sample weights are set to unity. If uncertainties are absent,
$\sigma_p=1$ defines the numerical metric and the residuals estimate its
scale. With neither input, minimizing $Q$ is equivalent to minimizing the
ordinary sample mean squared error. The solver applies linear least squares
to $W^{1/2}X$ and $W^{1/2}\bm y$ using \texttt{numpy.linalg.lstsq}; it
does not solve a nonlinear optimization problem. A rank-deficient selected
design matrix is rejected. For full column rank, the formal solution is
\begin{equation}
 \widehat{\bm c}=(X^{\trans}WX)^{-1}X^{\trans}W\bm y,
 \qquad \nu=N-K.
 \label{eq:solution}
\end{equation}
For equal-area Cartesian detector pixels, selecting exponent $s$ in the
standard analyzer also multiplies the fit weights by
\begin{equation}
 \rho_s(q)=
 \frac{(q_{\max}^2-q_{\min}^2)/2}
 {[q_{\max}^{s+1}-q_{\min}^{s+1}]/(s+1)}q^{s-1}.
 \label{eq:radial-fit-weight}
\end{equation}
The constant normalizes the area-average of $\rho_s$ to one and has no effect
on the fitted coefficients. Thus the basis orthonormality and detector
residual metric use the same $q^s\dd q\dd\theta$ measure. Supplied
uncertainty weights, sample weights, and fractional coverage multiply this
factor. Arbitrary point lists do not by themselves define a quadrature
measure; lower-level callers must supply the appropriate sample weights.
The reported \texttt{mse} always equals
$N^{-1}\sum_p(y_p-\widehat y_p)^2$, regardless of the fitting metric.
The weighted sum of squares is $Q(\widehat{\bm c})$; the field
\texttt{reduced\_chi\_square} equals $Q/\nu$. The latter has a conventional
reduced-$\chi^2$ interpretation only under the corresponding noise model.

Let $L=(X^{\trans}WX)^{-1}X^{\trans}W$ denote the linear map from
observations to coefficients. With absolute uncertainties, the supplied
$\sigma_p$ are the actual standard deviations of independent pixel
measurements, so $\Sigma_{\mathrm{abs}}=\operatorname{diag}(\sigma_p^2)$.
Additional sample weights change the fitting metric, not these measurement
variances. For $\nu>0$, the implemented coefficient covariance is
\begin{equation}
 C=\begin{cases}
 L\Sigma_{\mathrm{abs}}L^{\trans},
   &\text{absolute pixel uncertainties},\\
 \gamma(X^{\trans}WX)^{-1},
   &\text{relative precision model},
 \end{cases}
 \qquad \gamma=Q(\widehat{\bm c})/\nu.
 \label{eq:covariance}
\end{equation}
The first case is the default when experimental uncertainties are supplied
and implements the diagonal-noise specialization of the general
fixed-design sampling covariance
\begin{equation}
 C_{\mathrm{sampling}}=(X^{\trans}WX)^{-1}
 X^{\trans}W\Sigma WX(X^{\trans}WX)^{-1}.
 \label{eq:sandwich}
\end{equation}
For $w_p=1$ and $\Sigma=\Sigma_{\mathrm{abs}}$, it reduces to the usual
inverse normal matrix. For fractional area weights it generally does not:
the middle factor contains $w_p^2/\sigma_p^2$, rather than
$w_p/\sigma_p^2$. Multiplying all geometric weights by a common positive
constant leaves both the coefficient estimates and this absolute covariance
unchanged. Statistical precision factors must be incorporated in the
supplied $\sigma_p$ themselves.

The covariance is evaluated using a reduced QR factorization
$W^{1/2}X=UR$, without forming an inverse normal matrix or an $N\times N$
observation matrix. In the absolute case,
\begin{equation*}
 L\Sigma_{\mathrm{abs}}^{1/2}
 =R^{-1}U^{\trans}\operatorname{diag}(\sqrt{w_p}),
 \qquad C=(L\Sigma_{\mathrm{abs}}^{1/2})
           (L\Sigma_{\mathrm{abs}}^{1/2})^{\trans}.
\end{equation*}
The relative case uses triangular solves for
$\gamma R^{-1}R^{-\trans}$.

Relative uncertainties can explicitly request residual scaling; fits
without supplied uncertainties also use this second case. It retains the
working model $\Sigma=\gamma W^{-1}$, in which the combined fit weights
represent relative statistical precision. It does not calibrate geometric
coverage weights with otherwise unspecified pixel-noise variance. Merely
rescaling residuals does not supply the missing noise model. The
implementation returns an undefined covariance when $\nu\leq0$.

Returned errors and significance thresholds are conditional on the
selected noise and linear-signal model. General non-diagonal sample
covariances are not accepted; the errors do not include uncertainty in the fitted center,
mode selection, or reduction-induced correlations. For deterministic
simulation data, residual-scaled errors describe finite-expansion
discrepancies under that model, rather than experimental counting noise.

\subsection{Parseval identities and angular powers}

For any coefficient vector in the selected orthonormal basis, including a
detector-fitted vector, continuous orthonormality gives the exact identity
\begin{equation}
 \mathcal{P}=\|\widehat I\|_{\Ann,s}^2
 =\int_{q_{\min}}^{q_{\max}}\int_0^{2\pi}
  \widehat I^2q^s\dd\theta\dd q
 =\sum_{(n,m)\in\Modes}\widehat c_{nm}^{\,2}.
 \label{eq:parseval}
\end{equation}
For the infinite complete expansion this is Parseval's identity for $I$;
for a finite continuous projection the omitted-mode norm is the sum of the
omitted squared projection coefficients. Equation~\eqref{eq:parseval}
is exact for the finite reconstructed function even when its coefficients
were estimated from irregular or masked data. It does not equate its
continuous norm to the measured weighted residual statistic or to the
norm of an unknown signal in unobserved regions.

The separate angular powers are
\begin{align}
 \mathcal{P}_n^c&=\sum_{m:(n,m)\in\Modes}\widehat c_{n,m}^{\,2},
 &&n=0,1,\ldots,\nonumber\\
 \mathcal{P}_n^s&=\sum_{m:(-n,m)\in\Modes}\widehat c_{-n,m}^{\,2},
 &&n=1,2,\ldots,\qquad \mathcal{P}_0^s=0,
 \label{eq:angular-power}\\
 \mathcal{P}&=\sum_{n\geq0}(\mathcal{P}_n^c+\mathcal{P}_n^s).
 \nonumber
\end{align}
These powers quantify squared-intensity content on the annulus. They are
neither integrated scattering cross sections nor magnetic energies, and
they discard coefficient signs. For example, with
$M_s=2\pi\Delta_s/(s+1)$ and $\Phi_{00}^{(s)}=1/\sqrt{M_s}$,
the linear integral in the selected measure is
$\int\!\!\int\widehat Iq^s\dd\theta\dd q
=\sqrt{M_s}\,\widehat c_{00}$ when the constant mode is present.
If $q$ has units $Q$ and intensity has units $U$, then $R_m$ has units
$Q^{-(s+1)/2}$, $c_{nm}$ has units $UQ^{(s+1)/2}$, and $\mathcal{P}$
has units $U^2Q^{s+1}$.
Comparisons of absolute modal powers therefore require consistent
coordinate units, intensity normalization, and annulus limits.

For a branch block $\widehat{\bm c}_n^b$ and its marginal covariance block
$C_n^b$, with $b\in\{c,s\}$, the plotted power uncertainty is computed as
\begin{equation}
 \delta\mathcal{P}_n^b=
 \left[2\tr\!\left((C_n^b)^2\right)
 +4(\widehat{\bm c}_n^b)^{\trans}C_n^b\widehat{\bm c}_n^b\right]^{1/2}.
 \label{eq:power-error}
\end{equation}
This is the square root of the Gaussian quadratic-form variance evaluated
with the fitted coefficient vector, including the term quadratic in
covariance. This substitution gives a plug-in uncertainty estimate rather
than an exact sampling standard deviation with a known true coefficient
vector. Powers are
not noise-bias corrected; squaring noisy coefficients produces positive
power even for a zero true block. In the standard overview both angular
power panels use
\begin{equation*}
 M=\max_{n,b}\mathcal{P}_n^b,\qquad
 p_n^b=\frac{\mathcal{P}_n^b}{M},\qquad
 \delta p_n^b=\frac{\delta\mathcal{P}_n^b}{M}.
\end{equation*}
The same $M$ is used for both branches and is held fixed in error
propagation. The displayed values are not fractions of $\mathcal{P}$.
The power-axis label $P_n^2/P_{\max}^2$ in the current implementation
denotes $\mathcal{P}_n^b/M$; it does not square that power ratio again.
On the logarithmic axes, lower error bars are shortened where necessary
to remain above the positive plotting floor; the returned uncertainties
retain Eq.~\eqref{eq:power-error}. The separate $(n,m)$ power maps instead
display $\widehat c_{nm}^{\,2}/\max_j\widehat c_j^{\,2}$, using the largest
individual product-mode power as their common scale.

\subsection{Radial coefficient spectra associated with angular modes}

Keeping each angular order separate while summing its radial modes gives
the spectra returned by \texttt{angular\_radial\_spectra}:
\begin{equation}
 I_n^c(q)=\sum_m\widehat c_{n,m}R_m^{(s)}(q),\quad n\geq0,
 \qquad
 I_n^s(q)=\sum_m\widehat c_{-n,m}R_m^{(s)}(q),\quad n\geq1.
 \label{eq:radial-spectra}
\end{equation}
The stored sine spectrum at order zero is identically zero. These are
signed coefficient functions with intensity units, and the reconstruction is
\begin{equation}
 \widehat I(q,\theta)=\frac{I_0^c(q)}{\sqrt{2\pi}}
 +\frac{1}{\sqrt\pi}\sum_{n\geq1}
 \left[I_n^c(q)\cos(n\theta)+I_n^s(q)\sin(n\theta)\right].
 \label{eq:spectral-reconstruction}
\end{equation}
Writing the unnormalized real Fourier series as
$\widehat I=a_0+\sum_{n\geq1}[a_n\cos(n\theta)+b_n\sin(n\theta)]$,
the Fourier amplitudes and azimuthal average are
\begin{equation}
 a_0(q)=\langle\widehat I(q,\theta)\rangle_\theta
 =\frac{I_0^c(q)}{\sqrt{2\pi}},\qquad
 a_n(q)=\frac{I_n^c(q)}{\sqrt\pi},\qquad
 b_n(q)=\frac{I_n^s(q)}{\sqrt\pi}\quad(n\geq1).
 \label{eq:fourier-amplitudes}
\end{equation}
Thus $I_0^c$ itself is not the azimuthally averaged intensity. The spectra
retain radial sign changes and relative harmonic amplitudes that cannot be
recovered from angular powers alone. Radial orthonormality connects the
two descriptions exactly:
\begin{equation}
 \int_{q_{\min}}^{q_{\max}}[I_n^b(q)]^2q^s\dd q=\mathcal{P}_n^b.
 \label{eq:spectrum-power}
\end{equation}
In the unnormalized Fourier convention just specified, the same identities
read
\begin{align*}
 \mathcal{P}_0^c&=2\pi\int_{q_{\min}}^{q_{\max}}a_0(q)^2q^s\dd q,\\
 \mathcal{P}_n^c&=\pi\int_{q_{\min}}^{q_{\max}}a_n(q)^2q^s\dd q,\qquad
 \mathcal{P}_n^s=\pi\int_{q_{\min}}^{q_{\max}}b_n(q)^2q^s\dd q
 \quad(n\geq1).
\end{align*}
These factors are already included in the orthonormal coefficient
convention of Eq.~\eqref{eq:angular-power}.

Let $\bm R_n^b(q)$ contain the radial modes present in the selected branch
block. Covariance propagation gives
\begin{equation}
 \operatorname{Var}[I_n^b(q)]
 =[\bm R_n^b(q)]^{\trans}C_n^b\bm R_n^b(q),
 \qquad
 \operatorname{Cov}[I_n^b(q),I_n^b(q')]
 =[\bm R_n^b(q)]^{\trans}C_n^b\bm R_n^b(q').
 \label{eq:spectrum-error}
\end{equation}
The standard analyzer evaluates the spectra at 300 equally spaced $q$
values and displays pointwise one-standard-deviation bands from the first
expression. This grid samples the fitted function; it introduces no new
independent observations. Simultaneous confidence bands are not calculated.

\subsection{Piecewise-linear radial reference analysis}

As a reference that does not use global radial polynomials, the same angular
modes can be combined with local linear hat functions. Let
\begin{equation}
 q_{\min}=q_0<q_1<\cdots<q_K=q_{\max}
\end{equation}
define $K$ nonoverlapping radial intervals. The nodal functions $B_k(q)$ obey
$B_k(q_l)=\delta_{kl}$, are linear on each adjacent interval, and vanish
elsewhere. The reference expansion is
\begin{equation}
 \widehat I_{\mathrm{PL}}(q,\theta)
 =\sum_n\sum_{k=0}^{K}a_{n,k}B_k(q)\mathcal{T}_n(\theta),
 \qquad
 I_{n,\mathrm{PL}}(q)=\sum_{k=0}^{K}a_{n,k}B_k(q).
 \label{eq:piecewise-expansion}
\end{equation}
Hence $a_{n,k}=I_{n,\mathrm{PL}}(q_k)$ is directly the fitted polygonal
spectrum at a knot. The knots may be uniform in $q$ or equally spaced in the
selected radial measure,
\begin{equation}
 q_k=\left[q_{\min}^{s+1}
 +\frac{k}{K}\left(q_{\max}^{s+1}-q_{\min}^{s+1}\right)\right]^{1/(s+1)}.
 \label{eq:piecewise-knots}
\end{equation}
The fit uses the same admitted pixels, angular orders, uncertainties,
quadrature weights, and optional pixel integration as the polynomial branch.
No smoothing penalty is applied.

The hat functions overlap and are not orthogonal. Their coefficients must
therefore not be assigned powers by summing $a_{n,k}^2$. Instead define
\begin{equation}
 (G_s)_{kl}=\int_{q_{\min}}^{q_{\max}}B_k(q)B_l(q)q^s\dd q.
 \label{eq:piecewise-gram}
\end{equation}
The matrix $G_s$ is symmetric and tridiagonal. The reference angular power is
\begin{equation}
 \mathcal{P}_{n,\mathrm{PL}}^b
 =(\bm a_n^b)^{\trans}G_s\bm a_n^b,
 \label{eq:piecewise-power}
\end{equation}
which uses the same continuous norm as Eq.~\eqref{eq:angular-power}. For a
coefficient covariance $C_n^b$, its plug-in Gaussian uncertainty is
\begin{equation}
 \delta\mathcal{P}_{n,\mathrm{PL}}^b
 =\left[
 2\tr(G_sC_n^bG_sC_n^b)
 +4(\bm a_n^b)^{\trans}G_sC_n^bG_s\bm a_n^b
 \right]^{1/2}.
 \label{eq:piecewise-power-error}
\end{equation}
Agreement between the polynomial and polygonal $I_n(q)$ and angular powers
tests whether an inferred radial feature is robust to the chosen radial
representation.

\subsection{Angular block statistics \texorpdfstring{$T_n$}{Tn}}

Angular powers measure amplitude, whereas \texttt{angular\_block\_tests}
tests whether an entire angular contribution is distinguishable from zero
under the fit covariance model. For each branch and order the null is
$H_{0,n}^b:\bm c_n^b=\bm0$, with every other fitted coefficient left free.
The statistic and block rank are
\begin{equation}
 T_n^b=(\widehat{\bm c}_n^b)^{\trans}(C_n^b)^+_{\tau}
             \widehat{\bm c}_n^b,\qquad
 r_n^b=\#\{j:s_j>\tau_n^b\},
 \label{eq:block-statistic}
\end{equation}
where $r_n^b$ is the numerical rank. For a covariance block of dimension
$k_n^b$, a single singular-value decomposition determines both this rank
and the numerical pseudoinverse:
\begin{align*}
 C_n^b&=U\operatorname{diag}(s_j)V^{\trans},\qquad
 \tau_n^b=k_n^b\epsilon_{\mathrm{mach}}s_{\max},\\
 (C_n^b)^+_{\tau}&=V\operatorname{diag}(d_j)U^{\trans},\qquad
 d_j=\begin{cases}
 s_j^{-1},&s_j>\tau_n^b,\\
 0,&s_j\leq\tau_n^b.
 \end{cases}
\end{align*}
Here $\epsilon_{\mathrm{mach}}$ is the machine precision of the
double-precision covariance array. The same retained-singular-value mask
is used for the rank and pseudoinverse; values exactly at the cutoff are
discarded. The latter is the Moore--Penrose pseudoinverse of the truncated
covariance block and reduces to the ordinary inverse when all directions
are retained. All correlations between the fitted radial coefficients
enter the covariance block; individual squared coefficient-to-error
ratios cannot in general replace Eq.~\eqref{eq:block-statistic}. The
quantities $T_n^b$ are dimensionless test statistics, not expansion
coefficients or Parseval powers.

With independent Gaussian errors of known absolute scale and the propagated
sampling covariance in Eq.~\eqref{eq:covariance}, the null reference
distribution is $\chi^2_{r_n^b}$, including for additional geometric fit
weights. When the common variance scale is estimated from the same fit
residuals, the reference is $T_n^b/r_n^b\sim F_{r_n^b,\nu}$ under the
relative precision model $\Sigma=\gamma W^{-1}$. This F reference is not
calibrated for arbitrary geometric weights with a different noise model.
Thus the critical value returned at confidence $\beta$ is
\begin{equation}
 T_{n,\mathrm{crit}}^b=
 \begin{cases}
 \chi^2_{r_n^b;\,\beta},&\text{absolute covariance scale},\\
 r_n^b F_{r_n^b,\nu;\,\beta},&\text{residual-estimated covariance scale}.
 \end{cases}
 \label{eq:block-threshold}
\end{equation}
The default is $\beta=0.95$; upper-tail probabilities are also returned.
In the residual-scaled case the code evaluates the upper tail of the
$F_{r_n^b,\nu}$ distribution at $T_n^b/r_n^b$; the stored statistic itself
is not divided by $r_n^b$. The standard analyzer selects this case when
\texttt{covariance\_scaled\_by\_residual} is true and passes the fit's
residual degrees of freedom to the block-test function.
Missing blocks have rank zero and undefined statistics, including the
nonexistent zero-order sine block. A nominal rejection concerns the specified
block conditional on the other included modes, without correcting for the
other blocks tested. A small amplitude can have a large $T_n^b$ if it is
precisely determined. The noise-model qualifications following
Eq.~\eqref{eq:covariance} apply to these tests.

For nearly singular blocks the statistic tests only the covariance
directions retained by the common cutoff. The same numerical rank
determines the reference distributions and thresholds above. An existing
block with zero numerical rank has a zero stored statistic and undefined
upper-tail probability and critical value; it provides no calibrated test.

The block-test result additionally provides an optional numerical Holm
step-down correction for one joint family $\mathcal H$ containing all positive-rank cosine and sine
blocks for this map, fixed annulus, and selected basis. The family size is
$N_{\mathrm{test}}=|\mathcal H|$; missing and rank-zero blocks are excluded.
For the full rectangle with $n_{\max}=25$ and positive block ranks,
$N_{\mathrm{test}}=51$, independent of the radial truncation. Both branches
are corrected together, including the isotropic block, rather than defining
a separate family for each panel.

Let $p_{(1)}\leq\cdots\leq p_{(N_{\mathrm{test}})}$ denote the ordered
uncorrected upper-tail probabilities. The adjusted probabilities, mapped
back to their original blocks after sorting, are
\begin{equation}
 \widetilde p_{(j)}
 =\min\!\left\{1,\ \max_{1\leq i\leq j}
              \left[(N_{\mathrm{test}}-i+1)p_{(i)}\right]\right\},
 \qquad j=1,\ldots,N_{\mathrm{test}}.
 \label{eq:holm-pvalues}
\end{equation}
The rejection rule is $\widetilde p_n^b\leq\alpha$, where
$\alpha=1-\beta=0.05$ by default. Unavailable probabilities in positive-rank
blocks count conservatively as one in the adjustment and remain undefined
without rejection in the output. Missing and rank-zero blocks are excluded.

Provided the marginal probabilities are valid, Holm controls the
probability of rejecting any true null in this family by $\alpha$ under
arbitrary dependence. It does not repair an incorrect covariance model,
account for data-dependent basis or annulus selection, or control a larger
family spanning several maps or analysis runs. The $T_n^b$ statistics,
numerical ranks, raw probabilities, and nominal critical values remain
unchanged.

These error probabilities concern zero-coefficient null hypotheses under
the specified noise and signal model. They do not quantify the frequency
of experimental artifacts or physical misidentifications. A centering or
detector artifact can produce a genuinely nonzero fitted block and hence
a small model-based p-value. Calibration simulations check the
implementation for their generating models; they do not establish that
those assumptions hold for the measured data.

The standard overview shows the $T_n^b$ values but deliberately omits
critical-value curves and rejection markers. The nominal critical values,
raw probabilities, and numerical Holm properties remain available
separately; they do not determine a physical identification criterion,
modify the modal powers and spectra, or select and refit a reduced model.

\subsection{Parity, mirror symmetry, and signal centering}

Point reflection in the detector plane is $\theta\mapsto\theta+\pi$.
Both trigonometric branches obey
$\mathcal{T}_n(\theta+\pi)=(-1)^{|n|}\mathcal{T}_n(\theta)$.
An exactly parity-even pattern therefore has only even angular orders.
The package reports
\begin{equation}
 f_{\mathrm{odd}}=
 \frac{\sum_{|n|\,\mathrm{odd},m}\widehat c_{nm}^{\,2}}{\mathcal{P}},
 \qquad f_{\mathrm{even}}=1-f_{\mathrm{odd}},
 \label{eq:parity-fractions}
\end{equation}
provided $\mathcal{P}>0$. For the fitted expansion, Parseval gives
\begin{align}
 D_{\mathrm{par}}^2
 &=\frac{1}{M_s}\int_{q_{\min}}^{q_{\max}}\int_0^{2\pi}
 [\widehat I(q,\theta)-\widehat I(q,\theta+\pi)]^2
 q^s\dd\theta\dd q
 =\frac{4\mathcal{P}f_{\mathrm{odd}}}{M_s},
 \label{eq:parity-mismatch}\\
 d_{\mathrm{par}}
 &=\frac{\|\widehat I(q,\theta)-\widehat I(q,\theta+\pi)\|_{\Ann,s}}
          {\|\widehat I\|_{\Ann,s}}
 =2\sqrt{f_{\mathrm{odd}}}.
 \label{eq:relative-parity}
\end{align}
The first measure has intensity-squared units; the second is a scale-free
relative RMS mismatch and contains the same information as $f_{\mathrm{odd}}$.
These are diagnostics of the fitted function. With \texttt{even\_only=True},
the odd power and parity mismatch vanish by construction. Mirror symmetry
about the $q_z$ axis is $\theta\mapsto-\theta$ and forbids all sine coefficients;
it is an additional assumption independent of parity.

For an observable expected to be parity even, a broad background of odd
angular powers can serve as an empirical reference for measurement noise,
asymmetries, and processing artifacts. With the same annulus and power
normalization, an even power at or below that level is not clearly
resolved above this empirical background. This is a descriptive robustness
criterion, not a calibrated probability threshold: mode variances, radial
profiles, and sensitivities can differ. Physical odd contributions must
first be excluded before interpreting odd powers as a background.

For a cross section known to be parity even, an incorrect origin can create
odd harmonics. The optional centering step estimates a common displacement
of the reciprocal-space origin from antipodal intensity differences.
For a two-dimensional affine detector grid, the fitted coordinate map is
\begin{equation}
 \begin{pmatrix}q_z\\q_y\end{pmatrix}
 =B\begin{pmatrix}r\\s\end{pmatrix}+\bm b,\qquad
 \bm\delta=B\begin{pmatrix}\delta_r\\\delta_s\end{pmatrix},\qquad
 \bm q'_p=\bm q_p-\bm\delta,
 \label{eq:center-map}
\end{equation}
where $r,s$ are row and column indices. The antipodal partner in the
original coordinates is $2\bm\delta-\bm q_p$. Bilinear interpolation gives
$\widetilde I(2\bm\delta-\bm q_p)=\sum_j a_j I_j$ and conditional variance
$\widetilde\sigma_p^2=\sum_j a_j^2\sigma_j^2$. All neighbors with nonzero
interpolation weight must be valid. The minimized objective is
\begin{equation}
 Q_{\mathrm{cen}}(\bm\delta)=
 \sum_{p\in\mathcal{H}(\bm\delta)}
 \frac{[I_p-\widetilde I(2\bm\delta-\bm q_p)]^2}
      {\sigma_p^2+\widetilde\sigma_p^2}.
 \label{eq:centering-objective}
\end{equation}
Here $\mathcal{H}(\bm\delta)$ selects valid partners within the corrected
annulus and one representative half-plane, avoiding explicit antipodal
double counting. The admissible set can change with the trial shift. The
implementation minimizes the unnormalized sum, records the before/after
objectives and pair counts, and uses a bounded Powell search. The standard
analyzer bounds each detector-index displacement by two pixels by default;
the standalone grid refiner defaults to one pixel. Multiple images can be
passed to the standalone refiner to minimize the sum of their objectives
for one common shift.

An optional second, modular refinement operates directly on the modal
representation. For every trial reciprocal-space shift $\bm\delta$, the data
are refitted on a separately chosen finite basis $\Modes_{\mathrm{cen}}$, and
the minimized objective is
\begin{equation}
 J_{\mathrm{odd}}(\bm\delta)=
 \frac{\displaystyle\sum_{\substack{(n,m)\in\Modes_{\mathrm{cen}}\\
                         |n|\ {\rm odd}}}
       \widehat c_{nm}(\bm\delta)^2}
      {\displaystyle\sum_{(n,m)\in\Modes_{\mathrm{cen}}}
       \widehat c_{nm}(\bm\delta)^2}.
 \label{eq:odd-mode-centering}
\end{equation}
The standalone function
\texttt{refine\_center\_by\_odd\_mode\_power} accepts an explicit mode list
and a bounded shift in reciprocal-space units. The standard analyzer exposes
it through \texttt{odd\_mode\_centering=True}; by default it uses the compact
full rectangle $|n|\leq5$, $0\leq m\leq8$ and limits the additional shift to
$0.25$ nominal pixels around the current origin. It can start from the
nominal origin or locally refine the pixelwise parity result. The fitting
basis, annulus, measure exponent, sample weights, before/after objective and
power, selected-point counts, number of evaluations, and recovered shift are
retained as diagnostics. The final scientific fit remains independent of
the centering basis.

For a one-dimensional point list, the center is estimated on a regular
$129\times129$ auxiliary grid spanning $[-q_{\max},q_{\max}]^2$ by default.
Valid source points within $1.05q_{\max}$ are used for linear intensity
interpolation; centering uncertainties are assigned from the nearest
measured point. After refining this grid, the displacement is applied to
the original point coordinates. The final modal fit retains the original
intensities, uncertainties, and point sampling. In either centering path,
the selected annulus is reevaluated after the shift. Centering requires
supplied uncertainties and is disabled by default in the standard analyzer.
Geometry-dependent externally supplied weights must be recomputed for the
shifted origin; internal pixel integration computes coverage after centering.

Interpolation correlations and center-estimation uncertainty are not
propagated. Accordingly, neither $Q_{\mathrm{cen}}$ nor
$J_{\mathrm{odd}}$ is a calibrated goodness-of-fit statistic for the original
data. The modal optimum is additionally conditional on
$\Modes_{\mathrm{cen}}$, the annulus, weights, and search bounds. A physical
odd contribution must not be absorbed into a parity-based center.
In particular, the polarization-averaged spin-flip cross section and an
individual spin-flip channel need not have the same parity because an
individual channel can contain a chiral contribution. An odd-power fraction
alone therefore does not identify a reduction or centering error.

\subsection{Finite pixels, rebinning, and explicit preprocessing}

For an affine detector grid, optional subpixel quadrature estimates the
annular coverage of a pixel $P_p$ and the admitted-area average of each mode:
\begin{equation}
 f_p=\frac{|P_p\cap\Ann|}{|P_p|}
 \simeq\frac{N_{p,\mathrm{in}}}{s_{\mathrm{sub}}^2},\qquad
 \overline X_{pj}\simeq
 \frac{1}{N_{p,\mathrm{in}}}
 \sum_{\ell\in P_p\cap\Ann}\Phi_j(q_\ell,\theta_\ell).
 \label{eq:pixel-integration}
\end{equation}
Subpixel nodes lie at the centers of a regular
$s_{\mathrm{sub}}\times s_{\mathrm{sub}}$ subdivision. The affine map
accounts for pixel shape and orientation. There are two distinct uses:
\texttt{fractional\_annulus\_weights} supplies $f_p$ as a fit weight, while
\texttt{pixel\_integrated\_annulus\_design} also replaces the design row by
$\overline X_{pj}$. The latter admits any pixel with $f_p>0$ and multiplies
an existing sample weight by $f_p$ for $s=1$. For general $s$, its returned
quadrature weight is
\begin{equation}
 w_{p,s}\simeq\frac{1}{s_{\mathrm{sub}}^2}
 \sum_{\ell\in P_p\cap\Ann}\rho_s(q_\ell),
 \label{eq:pixel-measure-weight}
\end{equation}
which combines fractional coverage and the radial factor in
Eq.~\eqref{eq:radial-fit-weight}. The default standard analyzer instead
uses center-point selection and evaluation; supplying only external
fractional weights retains that center-point annulus selection. One must
therefore enable internal pixel integration to include pixels whose centers
lie outside the annulus but whose areas overlap it. Supplying the same
coverage externally as well would apply it twice.

Pixel averaging refers to the portion inside the analysis annulus; the
observed pixel intensity itself is not subdivided. This is an approximation
for boundary pixels that represent full-pixel averages. Pixel integration
also does not model wavelength spread or instrumental $q$ resolution.
With absolute errors, the supplied pixel variances remain unchanged by
coverage; Eq.~\eqref{eq:covariance} propagates them using the integrated
design and the effective fit weights.
For a point list without supplied pixel geometry, these helpers reduce to
binary annulus selection and point evaluation. The integrated design is
assembled in chunks to bound temporary storage.

The optional \texttt{cartesian\_bin\_average} routine converts point
samples into equal-area Cartesian cells by
\begin{equation}
 I_P=\frac{\sum_{\ell\in P}v_\ell I_\ell}
            {\sum_{\ell\in P}v_\ell}.
 \label{eq:rebinning}
\end{equation}
For a source grid uniform in $(q,\theta)$, $v_\ell=q_\ell$ supplies the
Jacobian of $q\dd q\dd\theta$; general nonuniform grids require their
actual quadrature weights. A duplicated angular endpoint at $2\pi$ should
be removed before averaging, as in the documented NuMagSANS data
preparation. Empty cells remain undefined. The routine returns occupancies
and source-weight sums, but does not propagate measurement uncertainties.
Equal-area Cartesian samples require no additional factor of $q$ in the
subsequent area-based fit.

Further preprocessing functions are explicit optional operations.
\texttt{symmetry\_diagnostic} compares original pixels to bilinearly
interpolated parity or mirror partners using residuals standardized by
$\sqrt{\sigma_p^2+\widetilde\sigma_p^2}$. Its summed statistic uses one
representative half-plane and is conditional on interpolation assumptions.
\texttt{symmetry\_fold} averages a complete symmetry orbit $\mathcal{O}_p$
by inverse variance,
\begin{equation}
 I_p^{\mathrm{fold}}=
 \frac{\sum_{j\in\mathcal{O}_p}I_j/\sigma_j^2}
      {\sum_{j\in\mathcal{O}_p}1/\sigma_j^2},\qquad
 (\sigma_p^{\mathrm{fold}})^2=
 \left(\sum_{j\in\mathcal{O}_p}1/\sigma_j^2\right)^{-1}.
 \label{eq:folding}
\end{equation}
Incomplete orbits are marked invalid. Parity folding imposes even orders;
adding mirror folding also imposes the absence of sine contributions.
The returned per-pixel errors omit correlations from interpolation and
shared orbit members. Folding is intended for visualization and sensitivity
checks; symmetry-constrained fits to the original samples avoid creating
these duplicated observations. Folding is never automatic.

\texttt{mask\_symmetry\_diagnostic} reports a weight-centroid displacement
normalized by the largest admitted radius, point-reflection residuals of
the mask, and the number of partially covered pixels.
\texttt{rotate\_coordinates} changes the angular origin according to
$\theta'=\theta-\varphi$ without changing intensities.
\texttt{rescale\_intensity} divides intensities and uncertainties by a
common positive number. \texttt{standardized\_residual\_outliers} flags
$|I_p-\widehat I_p|/\sigma_p>5$ by default and does not delete samples.
These functions do not implement experimental background, efficiency, or
polarization-leakage corrections; those belong to the preceding reduction.

\subsection{Implementation correspondence and data conventions}

The definitions above correspond to \path{basis.py} and \path{indices.py}
(basis and selected modes), \path{fit.py} (linear fit and covariance),
\path{power.py} (powers and parity), \path{spectra.py} (radial spectra), and
\path{block_test.py} (angular block tests). All source files lie under
\path{src/sans2d_mode_analyzer/}. Explicit detector operations reside in
\path{preconditioning.py}. The optional polygonal reference is implemented in
\texttt{piecewise\_linear\_}\allowbreak\texttt{basis.py} and
\texttt{piecewise\_linear\_}\allowbreak\texttt{reference.py}.
The \path{standard_analyzer} combines sample
selection, fitting over the selected modes, and all numerical diagnostics.
Its optional \texttt{radial\_measure\_}\allowbreak\texttt{exponent} applies
the same exponent to
the radial basis, detector-fit quadrature weights, Parseval diagnostics,
parity normalization, and reported radial spectra; the default is $s=1$.
It returns the fit, detector-space arrays, powers, block tests, parity
measures, and radial spectra independently of figure generation. The optional
overview is rendered separately by \path{figure_backend.py}; setting
\texttt{export\_figure=False} suppresses this rendering without changing the
numerical analysis.  Its final two rows show the signed spectra
$I_n^{c,s}(q)$ on linear axes and the corresponding magnitudes
$|I_n^{c,s}(q)|$ on common log-log axes.  The latter representation changes
only the visualization: it exposes weak contributions over several decades
but discards coefficient signs and turns every zero crossing into a downward
cusp.  No additional fit or power calculation is performed.

The ILL NeXus loader maps the stored \texttt{qx1} direction to $q_z$,
\texttt{qy1} to $q_y$, and converts the stored inverse-angstrom coordinates
to inverse nanometers by a default factor of ten. For the bundled
nanoflower data, the local data documentation identifies the 10 and
30~mT maps as $\tfrac12(I^{+-}+I^{-+})$, whereas the 100~mT and 1~T maps
contain the $+-$ channel alone. Their common baseline analysis uses
$0.015\leq q\leq0.100\,\mathrm{nm}^{-1}$, the supplied uncertainties,
$n_{\max}=m_{\max}=25$, $7\times7$ subpixel coverage weights, and no
parity centering or folding. Their supplied pixel errors are treated as
absolute, with covariance propagated through the coverage-weighted fit.
The ASCII loader reads columns $(q_z,q_y,I,\sigma)$ without converting
units; the legacy nanoporous-Fe benchmark has an unresolved coordinate-unit
caveat in its data README. Its numerical $q$ interval must therefore not
be assigned physical units without establishing that provenance.

A reproducible analysis specifies the input channel and reduction,
coordinate units and angular origin, annulus, retained mode set, weights
and covariance scaling, any centering or symmetry operations, the
truncation settings, and the test confidence and family.
The modal description then provides
a common representation for measured and simulated detector patterns,
with physical interpretation tied to the observable and assumptions used.

\section{Verification and interpretation strategy}
\label{sec:verification-strategy}

\subsection{White-Noise Test Model}

This test isolates a particularly simple question: if the true intensity
contains only a small number of occupied modes, how much fitted power should
appear in all other modes solely because independent white noise was added to
the detector pixels?  Answering this question provides a controlled reference
for the ``noise carpet'' in a modal power map.  It is a verification of the
linear fit and its covariance propagation, not a general noise model for
reduced experimental SANS data.

\paragraph{Synthetic data and repeated observations.}
Let $\bm c$ denote the known coefficient vector used to generate a synthetic
detector image and let $X$ be the same design matrix as in
Eq.~\eqref{eq:solution}.  One noisy observation is
\begin{equation}
 \bm y=X\bm c+\bm\epsilon,
 \qquad
 \bm\epsilon\sim\mathcal{N}(\bm 0,\sigma^2\mathbf{I}_N),
 \label{eq:white-noise-model}
\end{equation}
where every admitted detector pixel has the same noise variance and different
pixels are statistically independent.  Repeating the simulation means keeping
$X$, $\bm c$, and $\sigma$ fixed while drawing a new $\bm\epsilon$ for every
realization.  The modes used to generate the exact image are called
\emph{occupied}; all remaining fitted modes have a true coefficient equal to
zero and are called \emph{unoccupied}.  The latter are the positions at which
the noise floor is examined.

For ordinary least squares, the fitted coefficients satisfy
\begin{align}
 \widehat{\bm c}
 &=\bm c+A\bm\epsilon,
 &A&=(X^{\trans}X)^{-1}X^{\trans},
 \label{eq:white-noise-response}\\
 C=\operatorname{Cov}(\widehat{\bm c})
 &=\sigma^2AA^{\trans}
 =\sigma^2(X^{\trans}X)^{-1}.
 \label{eq:white-noise-covariance}
\end{align}
Thus the covariance is not obtained from the Monte Carlo calculation: it is
known analytically from the imposed pixel noise and the sampled basis.  The
Monte Carlo calculation independently checks whether repeated numerical fits
follow this prediction.  Equation~\eqref{eq:white-noise-covariance} is the
equal-variance specialization of Eq.~\eqref{eq:covariance}.

\paragraph{Why a zero coefficient has nonzero fitted power.}
The index $k$ below denotes one flattened product-mode index $(n,m)$.  Its
fitted coefficient is normally distributed,
\begin{equation}
 \widehat c_k\sim\mathcal{N}(c_k,C_{kk}).
 \label{eq:white-noise-coefficient}
\end{equation}
The displayed power of that product mode is the square
$p_k=\widehat c_k^2$.  Squaring removes the sign of a noise fluctuation and
therefore produces a positive average even when the true coefficient is zero:
\begin{align}
 \operatorname{E}[p_k]&=c_k^2+C_{kk},
 \label{eq:white-noise-power-mean}\\
 \operatorname{Var}(p_k)&=2C_{kk}^2+4c_k^2C_{kk}.
 \label{eq:white-noise-power-variance}
\end{align}
For an unoccupied mode, $c_k=0$, and these expressions simplify to
\begin{equation}
 \operatorname{E}[p_k]=C_{kk},\qquad
 \operatorname{Var}(p_k)=2C_{kk}^2,\qquad
 \frac{p_k}{C_{kk}}\sim\chi_1^2.
 \label{eq:white-noise-zero-mode}
\end{equation}
The diagonal covariance entry $C_{kk}$ is therefore the predicted mean noise
power at position $(n,m)$.  It is also the quantity labelled ``analytic mean''
in the Stage-2 power map.  The ``Monte Carlo mean'' is instead the arithmetic
mean of $\widehat c_k^2$ over many independently generated detector images.
Agreement between these two maps verifies the expected transfer of white
pixel noise into modal power.

Equation~\eqref{eq:white-noise-zero-mode} also explains why a single noisy
fit does not look like a perfectly flat carpet.  For one unoccupied mode the
standard deviation of $p_k$ is $\sqrt{2}$ times its mean.  Individual squared
coefficients consequently fluctuate strongly, even though their ensemble
mean is known exactly.  After averaging $R$ independent realizations, the
relative standard error of that mean is $\sqrt{2/R}$.  For $R=1000$ this is
about $4.5\%$ for each individual mode.

\paragraph{When is the carpet constant?}
White pixel noise does not by itself imply identical coefficient variances.
The expected floor is constant only when the diagonal entries $C_{kk}$ are
approximately equal.  Finite Cartesian sampling, an annular mask, missing
pixels, nonuniform uncertainties, or an ill-conditioned set of sampled modes
can make these entries different.  In the current Stage-2 test the predicted
zero-mode powers vary by only about $2\%$ around their mean, so describing the
ensemble-averaged result as an approximately constant carpet is appropriate.
This near constancy is a property of this design matrix, rather than a general
theorem about white noise.

\paragraph{Powers summed over several radial modes.}
The angular powers in Eq.~\eqref{eq:angular-power} sum several squared
coefficients.  For a block $\mathcal{B}$ with coefficient subvector
$\widehat{\bm c}_{\mathcal{B}}$ and covariance block $C_{\mathcal{B}}$, the
corresponding power
$p_{\mathcal{B}}=\widehat{\bm c}_{\mathcal{B}}^{\trans}
\widehat{\bm c}_{\mathcal{B}}$ obeys
\begin{align}
 \operatorname{E}[p_{\mathcal{B}}]
 &=\bm c_{\mathcal{B}}^{\trans}\bm c_{\mathcal{B}}
   +\tr(C_{\mathcal{B}}),
 \label{eq:white-noise-block-mean}\\
 \operatorname{Var}(p_{\mathcal{B}})
 &=2\tr(C_{\mathcal{B}}^2)
   +4\bm c_{\mathcal{B}}^{\trans}C_{\mathcal{B}}\bm c_{\mathcal{B}}.
 \label{eq:white-noise-block-variance}
\end{align}
For a true zero block, its mean noise power is therefore
$\tr(C_{\mathcal{B}})$ and its variance is
$2\tr(C_{\mathcal{B}}^2)$.  Coefficient correlations enter through the
latter trace.  In the special case of $M$ independent coefficients with a
common variance $v$, one has
$p_{\mathcal{B}}/v\sim\chi_M^2$ and a relative standard deviation
$\sqrt{2/M}$.  Summed angular powers can consequently appear substantially
smoother than the powers of individual $(n,m)$ modes.

\paragraph{Numerical realization of the test.}
The reproducible script
\path{verification_examples/stage_2/example_1.py} uses the deterministic
Stage-1 signal, containing four occupied even-angular modes, and fits the full
rectangle $|n|\leq8$, $0\leq m\leq8$.  It adds Gaussian pixel noise with
$\sigma$ equal to $3\%$ of the maximum absolute exact intensity and evaluates
1000 independent noise realizations.  In the resulting $(n,m)$ maps, positive
$n$ denotes cosine modes, negative $n$ denotes sine modes, and the four white
cells mark the occupied generating modes that are deliberately excluded from
the zero-mode comparison.  The current test gives an empirical-to-analytic
mean-power ratio of $1.0010$ and an empirical-to-analytic power-variance ratio
of $1.0017$.  The normalized powers of the unoccupied modes also follow the
$\chi_1^2$ density in Eq.~\eqref{eq:white-noise-zero-mode}.  These agreements
verify the coefficient covariance and the positive noise bias of squared
modal amplitudes under the stated assumptions.

No Monte Carlo calculation is required to obtain Eqs.~\eqref{eq:white-noise-power-mean}--
\eqref{eq:white-noise-block-variance}; it is used here as a transparent
numerical verification.  Experimental data can contain counting-statistical
heteroscedasticity, background-subtraction errors, interpolation artifacts,
pixel-to-pixel correlations, centering uncertainty, and model truncation
error.  In those cases the white-noise test remains a useful regression test,
but its analytic carpet must not be interpreted as a calibrated experimental
significance threshold unless the corresponding experimental covariance
model has independently been justified.

\subsection{Heteroskedastic Gaussian Test Model}

The Stage-3 test removes the equal-variance assumption while retaining
independent Gaussian detector errors.  Its purpose is to verify that known,
pixel-dependent experimental uncertainties are propagated through the
weighted fit correctly.  The observation model is
\begin{equation}
 \bm y=X\bm c+\bm\epsilon,
 \qquad
 \bm\epsilon\sim\mathcal{N}(\bm0,\Sigma),
 \qquad
 \Sigma=\operatorname{diag}(\sigma_1^2,\ldots,\sigma_N^2).
 \label{eq:heteroskedastic-model}
\end{equation}
The errors remain statistically independent, but their variances need not be
equal.  This distinction is important: heteroskedastic noise is not white
noise, even though both models have a diagonal pixel covariance.

\paragraph{Weighted estimator and analytic covariance.}
When the supplied standard deviations are the true measurement errors and no
additional sample weights are used, inverse-variance weighting gives
\begin{align}
 W&=\Sigma^{-1},\nonumber\\
 \widehat{\bm c}
 &=\bm c+L\bm\epsilon,
 &L&=(X^{\trans}\Sigma^{-1}X)^{-1}X^{\trans}\Sigma^{-1},
 \label{eq:heteroskedastic-response}\\
 C=\operatorname{Cov}(\widehat{\bm c})
 &=L\Sigma L^{\trans}
 =(X^{\trans}\Sigma^{-1}X)^{-1}.
 \label{eq:heteroskedastic-covariance}
\end{align}
This is the inverse-variance specialization of the sandwich covariance in
Eq.~\eqref{eq:sandwich}.  The final equality in
Eq.~\eqref{eq:heteroskedastic-covariance} relies on using the correct
$W=\Sigma^{-1}$.  If the supplied uncertainties do not describe the actual
pixel variances, the fit remains numerically defined but its reported
covariance and probabilistic calibration need not be correct.

Because Eq.~\eqref{eq:heteroskedastic-model} is Gaussian and the estimator is
linear,
\begin{equation}
 \widehat c_k\sim\mathcal{N}(c_k,C_{kk}),
 \qquad
 c_k=0\ Longrightarrow\
 \frac{\widehat c_k^2}{C_{kk}}\sim\chi_1^2.
 \label{eq:heteroskedastic-zero-mode}
\end{equation}
Thus the normalized zero-mode power has exactly the same reference
distribution as in Eq.~\eqref{eq:white-noise-zero-mode}.  What changes is the
unnormalized mean power $C_{kk}$: unequal pixel errors alter the fitted-mode
covariance through the complete sampled basis and generally produce a
mode-dependent noise floor.  Dividing every mode by one common empirical
floor would therefore discard relevant information.  The appropriate
normalization is the individual analytic variance $C_{kk}$.

\paragraph{Controlled radial uncertainty profile.}
To make this effect visible without introducing a physical SANS count model,
the Stage-3 example prescribes a simple radial standard-deviation profile.
For an admitted pixel $p$, define
\begin{equation}
 u_p=\frac{q_p-q_{\min}}{q_{\max}-q_{\min}},
 \qquad
 \sigma_p=0.015 I_{\max}(1+3u_p),
 \qquad
 I_{\max}=\max_p|I_p^{\mathrm{exact}}|.
 \label{eq:heteroskedastic-profile}
\end{equation}
The pixel standard deviation therefore rises from $1.5\%$ to $6\%$ of the
maximum exact intensity between the inner and outer annulus boundaries.  This
profile is deliberately chosen as a transparent numerical test; it is not
claimed to be a detector or counting-statistical model.

The script \path{verification_examples/stage_3/example_1.py} uses the same
four-mode exact signal and full $|n|\leq8$, $0\leq m\leq8$ fit as Stage~2.
It propagates Eq.~\eqref{eq:heteroskedastic-profile} analytically and compares
the result with 1000 independent noise realizations.  The expected powers of
the 149 unoccupied modes have a relative variation of about $10.0\%$, compared
with about $2\%$ in the equal-variance Stage-2 design.  Nevertheless, the
empirical-to-analytic mean-power ratio is $0.9982$, the corresponding
power-variance ratio is $0.9943$, and the normalized powers follow the
$\chi_1^2$ reference distribution.  The structured raw power floor is
therefore predicted by the weighted design rather than indicating occupied
physical modes.

\paragraph{Regression criterion.}
The automated test \path{tests/test_stage_3_regression.py} generates 1024
independent realizations with a fixed random seed and passes the known
$\sigma_p$ values to the absolute-uncertainty fit.  It verifies the reported
coefficient covariance against Eq.~\eqref{eq:heteroskedastic-covariance}.  For
each unoccupied mode it then forms
\begin{equation}
 z_{rk}=\frac{\widehat c_{rk}-c_k}{\sqrt{C_{kk}}},
 \label{eq:heteroskedastic-standardized-error}
\end{equation}
where $r$ labels the realization, and checks
$\operatorname{E}[z_k]=0$, $\operatorname{E}[z_k^2]=1$,
$\operatorname{Var}(z_k^2)=2$, and $95\%$ coverage of the
$\chi_1^2$ quantile.  The acceptance tolerances are five analytic Monte Carlo
standard errors for 1024 realizations.  They are therefore tied to the finite
ensemble size rather than to the particular values generated by the fixed
seed.  This regression test verifies the complete path from a supplied
heteroskedastic error map through weighted fitting to calibrated modal power;
it does not establish that an experimental error map is itself correct.

\subsection{Correlated Gaussian Test Model}

Stage~4 retains Gaussianity but allows different detector pixels to fluctuate
coherently.  This is the first verification stage for which a map of marginal
pixel errors is not sufficient to determine the coefficient covariance.  A
computationally convenient low-rank model is
\begin{equation}
 \bm\epsilon=D\bm z+H\bm a,
 \qquad
 \bm z\sim\mathcal{N}(\bm0,\mathbf I_N),
 \qquad
 \bm a\sim\mathcal{N}(\bm0,\mathbf I_R),
 \qquad
 \Sigma=DD^{\trans}+HH^{\trans},
 \label{eq:correlated-model}
\end{equation}
where $D$ is diagonal and the $R$ columns of $H$ are fixed spatial
background patterns with independent random amplitudes.  The two Gaussian
vectors are mutually independent.  This factorized construction permits
sampling and covariance propagation without forming a dense $N\times N$
pixel covariance.

The marginal error reported for pixel $p$ is
\begin{equation}
 \sigma_{p,\mathrm{marg}}^2=\Sigma_{pp}
 =D_{pp}^2+\sum_{r=1}^{R}H_{pr}^2.
 \label{eq:correlated-marginal}
\end{equation}
These numbers describe the diagonal of $\Sigma$ but contain no information
about its off-diagonal entries.  Two noise models can therefore have identical
error bars at every pixel and nevertheless produce very different fitted-mode
uncertainties.

\paragraph{True and diagonal coefficient covariance.}
The current standard analyzer accepts independent pixel uncertainties rather
than a general pixel covariance.  When it is supplied only with the marginal
errors in Eq.~\eqref{eq:correlated-marginal}, it uses
\begin{equation}
 \Sigma_{\mathrm{diag}}=\operatorname{diag}(\Sigma),
 \qquad W=\Sigma_{\mathrm{diag}}^{-1},
 \qquad
 L_d=(X^{\trans}WX)^{-1}X^{\trans}W.
 \label{eq:correlated-diagonal-response}
\end{equation}
The resulting coefficient estimate remains linear and unbiased for a
zero-mean error, but two distinct covariance matrices must be distinguished:
\begin{align}
 C_{\mathrm{reported}}
 &=L_d\Sigma_{\mathrm{diag}}L_d^{\trans},
 \label{eq:correlated-reported-covariance}\\
 C_{\mathrm{true}}
 &=L_d\Sigma L_d^{\trans}
 =L_d(DD^{\trans}+HH^{\trans})L_d^{\trans}.
 \label{eq:correlated-true-covariance}
\end{align}
Only $C_{\mathrm{true}}$ describes repeated measurements generated by
Eq.~\eqref{eq:correlated-model}.  For a true zero coefficient,
\begin{equation}
 \frac{\widehat c_k^2}{(C_{\mathrm{true}})_{kk}}\sim\chi_1^2,
 \qquad
 \frac{\widehat c_k^2}{(C_{\mathrm{reported}})_{kk}}
 \sim r_k\chi_1^2,
 \qquad
 r_k=\frac{(C_{\mathrm{true}})_{kk}}
 {(C_{\mathrm{reported}})_{kk}}.
 \label{eq:correlated-power-scaling}
\end{equation}
Consequently, a diagonal error model can produce either understated or
overstated mode significance.  In particular, a smooth correlated pattern
can project almost entirely onto one low angular order, even though its
marginal pixel variance appears modest.

If the full positive-definite $\Sigma$ were known, the efficient
correlation-aware estimator would instead be generalized least squares,
\begin{equation}
 \widehat{\bm c}_{\mathrm{GLS}}
 =(X^{\trans}\Sigma^{-1}X)^{-1}X^{\trans}\Sigma^{-1}\bm y,
 \qquad
 C_{\mathrm{GLS}}=(X^{\trans}\Sigma^{-1}X)^{-1}.
 \label{eq:correlated-gls}
\end{equation}
Stage~4 does not add a dense-covariance GLS interface to the standard
analyzer.  It characterizes the failure of the existing diagonal statistical
model when the data-generating process is correlated.

\paragraph{Controlled correlated background.}
The script \path{verification_examples/stage_4/example_1.py} uses the same
exact signal and modal rectangle as Stages~2 and 3.  Its independent component
has $D_{pp}=0.020I_{\max}$.  The two columns of $H$ are
\begin{equation}
 h_{1,p}=0.010I_{\max},
 \qquad
 h_{2,p}=0.030I_{\max}\frac{q_{y,p}}{q_{\max}}.
 \label{eq:correlated-patterns}
\end{equation}
The first is a common offset and the second is a smooth vertical gradient.
The gradient is parity odd and therefore couples strongly to sine modes of odd
angular order.  These deliberately simple patterns are diagnostic test cases,
not a proposed instrumental model for a particular experiment.

For 1000 realizations, the empirical-to-true mean-power ratio of the 149
unoccupied modes is $0.9973$ and the corresponding power-variance ratio is
$1.0018$.  Normalization by $C_{\mathrm{true}}$ therefore reproduces the
analytic Gaussian result.  In contrast, for the most affected mode
$(n,m)=(-1,0)$ the true coefficient variance is $674.2$ times the variance
reported by the diagonal model.  Its nominal $95\%$ coverage is $94.9\%$ with
the true covariance but only $5.9\%$ with the diagonal covariance.  The
resulting apparently significant odd mode is generated entirely by the
correlated background.

The automated test \path{tests/test_stage_4_regression.py} repeats this
experiment with a smaller design and 1024 realizations.  It verifies the
$\chi_1^2$ moments and coverage under $C_{\mathrm{true}}$, confirms that the
public fit reports Eq.~\eqref{eq:correlated-reported-covariance}, and requires
the correlation-induced miscalibration to remain detectable.  This is a
regression characterization of a known model limitation, rather than a claim
that marginal experimental error bars determine the required full covariance.

\subsection{Detector-Centre Translation Test}

Stage~5 isolates a deterministic geometrical systematic rather than another
noise distribution.  The synthetic generator accepts
\texttt{center\_shift=(delta\_q\_z, delta\_q\_y)} and evaluates the prescribed
modal signal about that displaced physical centre while retaining laboratory
detector coordinates.  This makes it possible to generate a signal whose
true expansion contains only even angular orders, analyze it about an
incorrect nominal origin, and then repeat the same analysis after the
parity-based centre refinement of Eq.~\eqref{eq:centering-objective}.  No
random noise is required for this test.  Unit uncertainties merely define
equal numerical weights and do not represent a statistical error model.

\paragraph{First-order angular-mode coupling.}
Let the parity-even signal contain only the two unnormalized Fourier
components
\begin{equation}
 I_{\mathrm{even}}(q,\theta)=A_0(q)+A_2(q)\cos(2\theta).
 \label{eq:stage5-even-signal}
\end{equation}
If its physical centre is at
$\bm\delta=(\delta_z,\delta_y)$ but the pattern is analyzed about the origin,
the observed detector function is
\begin{equation}
 I_{\bm\delta}(q_z,q_y)
 =I_{\mathrm{even}}(q_z-\delta_z,q_y-\delta_y)
 =I_{\mathrm{even}}-\delta_z\partial_{q_z}I_{\mathrm{even}}
  -\delta_y\partial_{q_y}I_{\mathrm{even}}+O(|\bm\delta|^2).
 \label{eq:stage5-translation}
\end{equation}
Using
\begin{equation}
 \partial_{q_z}=\cos\theta\,\partial_q
 -\frac{\sin\theta}{q}\partial_\theta,
 \qquad
 \partial_{q_y}=\sin\theta\,\partial_q
 +\frac{\cos\theta}{q}\partial_\theta,
 \label{eq:stage5-polar-derivatives}
\end{equation}
the displacement along $q_y$ gives
\begin{align}
 \partial_{q_y}I_{\mathrm{even}}
 ={}&\left(A_0'-\frac12A_2'-\frac{A_2}{q}\right)\sin\theta
 \nonumber\\
 &+\left(\frac12A_2'-\frac{A_2}{q}\right)\sin(3\theta),
 \label{eq:stage5-qy-coupling}
\end{align}
where primes denote radial derivatives.  Thus a small vertical translation
of an exact $n=0,2$ cosine signal produces the distinctive odd sine orders
$n=1$ and $n=3$.  The corresponding $q_z$ derivative is
\begin{align}
 \partial_{q_z}I_{\mathrm{even}}
 ={}&\left(A_0'+\frac12A_2'+\frac{A_2}{q}\right)\cos\theta
 \nonumber\\
 &+\left(\frac12A_2'-\frac{A_2}{q}\right)\cos(3\theta),
 \label{eq:stage5-qz-coupling}
\end{align}
so a horizontal translation couples to the cosine branches of the same odd
orders.  A general displacement produces both patterns.  The minus signs in
Eq.~\eqref{eq:stage5-translation} determine the signs of the induced spectra.
These relations show why coherent radial profiles in $I_1^{c,s}(q)$ and
$I_3^{c,s}(q)$ can diagnose a centre displacement more specifically than a
single total odd-power fraction.

\paragraph{Controlled numerical realization.}
The script \path{verification_examples/stage_5/example_1.py} generates the
three exact product modes $(n,m)=(0,0),(0,2),(2,1)$ on
$0.05\leq q\leq0.30$ and displaces their common centre by $0.35$ nominal
pixels along $q_y$, or
$\delta_y=2.625\times10^{-3}$ in the chosen reciprocal-coordinate units.
The raw full-rectangle fit has an odd-power fraction
$f_{\mathrm{odd}}=6.9043\times10^{-4}$.  Its two dominant odd sine powers
occur at $n=1$ and $n=3$, and their fitted radial spectra agree with
Eq.~\eqref{eq:stage5-qy-coupling} to first order.  The pixelwise centre refinement
returns
$(\widehat\delta_z,\widehat\delta_y)
=(6.7530\times10^{-5},2.62555\times10^{-3})$ and reduces the fitted odd-power
fraction to $1.2879\times10^{-7}$.  It also recovers the three generating
even-mode coefficients to the displayed numerical precision.  The remaining
$q_z$ displacement is an interpolation-dependent numerical residual and
produces small odd cosine powers. Starting from this result, the additional
modal refinement of Eq.~\eqref{eq:odd-mode-centering} returns the total shift
$(5.2\times10^{-12},2.625000007\times10^{-3})$ and reduces the odd-power
fraction to $5.95\times10^{-21}$. The example writes separate uncorrected,
pixel-centred, and modal-refined standard overviews and a direct comparison of
the induced $I_1^s(q)$ and $I_3^s(q)$ spectra with their first-order
predictions.

The regression test \path{tests/test_stage_5_regression.py} fixes this
noise-free construction.  It requires recovery of both centre coordinates to
within $0.02$ nominal pixels, identifies $n=1,3$ as the two strongest raw odd
sine orders, and requires the centre correction to suppress the odd-power
fraction by more than three orders of magnitude while recovering the known
even coefficients. It additionally requires the bounded modal refinement to
recover the imposed shift to $10^{-8}$ in both reciprocal-space coordinates
and to reduce the final odd-power fraction below $10^{-16}$.

\paragraph{Interpretational boundary.}
This test demonstrates that a small origin error is sufficient to create the
kind of coherent odd-mode structure seen in an uncentered detector map.  If
the structure disappears after a small, independently plausible centre
correction and the recovered $n=1,3$ spectra follow the translation formulas,
a geometrical origin is strongly supported.  The inference remains
conditional on the physical parity assumption used by the centering
objective.  A genuinely parity-odd contribution can also be reduced by that
operation and must not be declared instrumental merely because it disappears.
Minimizing all odd modes is therefore not a more model-independent operation
than pixelwise parity centering; it is a useful numerical cross-check when
the parity constraint has already been established physically. Stability
against reasonable changes of $\Modes_{\mathrm{cen}}$, annulus, and shift
bound should be checked for measured data.
Moreover, the present fit treats the recovered centre as fixed; its estimation
uncertainty and the interpolation correlations are not propagated into the
reported modal covariance.  Stage~5 is therefore a systematic-error
verification and sensitivity test, not a calibrated hypothesis test that
excludes all physical odd contributions.



\end{document}
